% TOP_PROBLEM: {"id": 1424, "title": "Sumner's Conjecture", "category_id": 3, "category": {"id": 3, "name": "graph_theory", "display_name": "Graph Theory", "description": "Problems involving graphs, networks, and their properties.", "slug": "graph-theory", "order_index": 3, "created_at": "2026-07-31T15:26:25.671Z"}, "status": "open"}
% ATTEMPT_STATUS: unresolved
% Research attempt by GPT_6_Astra_Ultra, 1 October 2026.
% Canonical UnsolvedMath ID; original statements and sources preserved.
% FIRST_POSTED: 2026-10-01
% Imported from: unsolvedmath_additional_500.tex; UnsolvedMath ID 1424
% !TeX program = lualatex
% Compile twice: lualatex 1424.tex
% Self-contained: no external bibliography, images, or \input files.
% Numeric section labels and visible numbers are original UnsolvedMath IDs.
\documentclass[11pt,a4paper]{article}
\usepackage[margin=24mm,headheight=14pt]{geometry}
\usepackage{fontspec}
\setmainfont{Latin Modern Roman}
\setsansfont{Latin Modern Sans}
\setmonofont{Latin Modern Mono}
\usepackage{amsmath,amssymb,mathtools}
\usepackage{unicode-math}
\setmathfont{Latin Modern Math}
\usepackage{microtype}
\usepackage{enumitem,xurl,needspace}
\usepackage[dvipsnames]{xcolor}
\usepackage{hyperref}
\hypersetup{unicode=true,colorlinks=true,linkcolor=MidnightBlue,urlcolor=MidnightBlue,
 pdftitle={UnsolvedMath Problem 1424},pdfauthor={GPT 6 Astra Ultra},bookmarksnumbered=true}
\usepackage{bookmark,tocloft,titlesec,fancyhdr}
\setlength{\cftsecnumwidth}{4.2em}
\cftsetpnumwidth{2.5em}
\cftsetrmarg{4em}
\titleformat{\section}{\Large\bfseries\raggedright}{\thesection}{0.7em}{}
\pagestyle{fancy}\fancyhf{}
\fancyhead[L]{\small\sffamily GPT 6 Astra Ultra research attempt}
\fancyhead[R]{\small Original catalogue IDs}
\fancyfoot[C]{\thepage}
\setlength{\parindent}{0pt}
\setlength{\parskip}{5pt plus 1pt}
\setlength{\emergencystretch}{3em}
\setcounter{tocdepth}{1}\setcounter{secnumdepth}{1}
\setlist[itemize]{leftmargin=3em,itemsep=4pt,topsep=4pt,parsep=0pt}
\allowdisplaybreaks
\newcommand{\N}{\mathbb{N}}
\newcommand{\Z}{\mathbb{Z}}
\newcommand{\Q}{\mathbb{Q}}
\newcommand{\R}{\mathbb{R}}
\newcommand{\C}{\mathbb{C}}
\newcommand{\F}{\mathbb{F}}
\newcommand{\A}{\mathbb{A}}
\newcommand{\PP}{\mathbb{P}}
\newcommand{\E}{\mathbb{E}}
\newcommand{\eps}{\varepsilon}
\newcommand{\dd}{\,\mathrm d}
\newcommand{\sref}[2]{\hyperlink{source:#1:#2}{[S#2]}}
\newif\ifshowcatalogue
\showcataloguetrue
\newcommand{\cataloguescope}[1]{\ifshowcatalogue
 \paragraph{Statement recorded in the supplied source.}
 #1\fi}
\begin{document}
% UnsolvedMath ID 1424; source catalogue code GRAPH-037

\setcounter{section}{1423}

\section{Sumner’s Tournament Embedding Conjecture}\label{1424}

{\small\sffamily UnsolvedMath ID 1424 · Graph Theory}

\subsection{Definitions and mathematical statement}

\paragraph{Shared notation.}Unless a section says otherwise, $\N=\{1,2,\ldots\}$,
$\Z$ is the ring of integers, $\Q,\R,\C$ are the rational, real, and complex
fields, $[N]=\{1,\ldots,\lfloor N\rfloor\}$, and $\log$ is the natural
logarithm. For real functions of a parameter tending to infinity,
$f\ll g$ and $f=O(g)$ mean $|f|\leq Cg$ eventually for some fixed $C$;
$f\gg g$ reverses this comparison; $f=o(g)$ means $f/g\to0$;
$f\sim g$ means $f/g\to1$; and $f\asymp g$ means both order bounds.
A subscript on $O$ or $\ll$ specifies allowed constant dependence.
The expression $n^{o(1)}$ in an upper bound means growth smaller than
$n^\epsilon$ for each fixed $\epsilon>0$. Local definitions govern any
exception, finite-field convention, or use of the same symbol for another
object. An asymptotic claim is not automatically an effective algorithm.



A tournament chooses exactly one orientation for each edge of a complete graph. An oriented tree is an arbitrary orientation of a finite undirected tree. For each $n\geq2$, ask whether every tournament on $2n-2$ vertices contains every oriented tree on $n$ vertices as an orientation-preserving subgraph; induced embedding is not required. \sref{1424}{1} \sref{1424}{2}

\paragraph{Statement recorded in the supplied source.}
Does every (2n-2)-vertex tournament contain every n-vertex oriented tree? \sref{1424}{1}

\paragraph{Residual task reported by the archive.}

The embedded review (2026-08-17) records: Resolve the finite range not covered by the sufficiently-large-n theorem, or give a genuinely uniform proof of f(n)=2n-2. \sref{1424}{1}

\subsection{Short English statement}

Does every tournament on twice n minus two vertices contain every possible orientation of an n-vertex tree?

\subsection{Research attempt}

\paragraph{Uniform statement and known large-order theorem.}
The target is every $n\ge2$ and every orientation of every $n$-vertex
tree. The primary abstract of K\"uhn, Mycroft, and Osthus [E1] proves
the conjecture for sufficiently large $n$. Ai and Lian's 2026 abstract
[E2] gives the uniform upper bound
$f(n)\le\lceil(18n-23)/7\rceil$. These assertions have different
quantifiers and neither abstract by itself supplies the missing uniform
sharp result. This attempt proves the sharp obstruction, two complete
tree families, and the smallest orders directly.

\paragraph{The number $2n-2$ is necessary.}
For $n\ge2$, put $m=2n-3$ and take vertices in $\Z/m\Z$.
When $m>1$, orient the edge from $i$ to $j$ if
$j-i\pmod m\in\{1,\ldots,n-2\}$. Exactly one of the two
opposite nonzero differences is in this set, so this is a tournament,
and every vertex has outdegree $n-2$. It contains no out-star with
$n-1$ leaves, since the image of the centre would need outdegree
at least $n-1$. At $n=2$ the one-vertex tournament supplies the
same obstruction trivially. Thus any uniform unavoidable order satisfies
\[
                              f(n)\ge2n-2.
\]
This establishes sharpness of the proposed constant if the upper bound
is eventually proved. It does not establish that upper bound.

\paragraph{Both orientations of stars satisfy the exact bound.}
In any tournament on $2n-2$ vertices, the sum of all outdegrees is
$\binom{2n-2}{2}$. Their average is $(2n-3)/2=n-3/2$, so some
integer outdegree is at least $n-1$. Choose that vertex for the centre
and any $n-1$ of its outneighbours for the leaves. This embeds the
out-star. Applying the same calculation to indegrees embeds the
in-star. Edges between the leaves are irrelevant, since the embedding
is not required to be induced. Ignoring that last convention would
incorrectly rule out every nontrivial tree in a tournament.

\paragraph{Consistently directed paths require even fewer vertices.}
Every finite tournament has a directed Hamiltonian path. Prove this by
induction: take a directed path $v_1\to\cdots\to v_s$ through the
old vertices and insert a new vertex $w$. If $w\to v_1$, insert it
first. Otherwise find the first $i$ with $w\to v_i$; minimality
implies $v_{i-1}\to w$, so insertion between $v_{i-1}$ and $v_i$
works. If no such $i$ exists, every $v_i\to w$, and appending $w$
works. Every consecutive edge retains the required orientation.

Taking any $n$-vertex subtournament and this path embeds the oriented
tree whose edges all point along one path direction. This proof does
not cover arbitrary alternating orientations of a path: a directed
Hamiltonian path has one consistent orientation, and reversing selected
edges is not an available operation on the tournament. The distinction
matters when identifying the family actually proved.

\paragraph{All oriented trees of orders two and three.}
At order two, any tournament with two vertices has one oriented edge,
which embeds either labeling of the two-vertex oriented tree. At order
three the underlying tree is a path with two edges. Its orientation
is an in-star, an out-star, or a consistently directed path, up to
relabeling. The preceding proofs therefore settle every orientation at
$n=3$ in a four-vertex tournament, as required.

The executed check enumerates all $2^{\binom42}=64$ labeled tournaments
on four vertices and all four edge orientations of the labeled path
$0-1-2$. It tries all injective maps of its vertices into the
tournament and verifies an embedding in each of the $256$ cases.
The exhaustive check agrees with the symbolic classification; it is
not used to extrapolate to larger trees.

\paragraph{The obstruction to the most immediate leaf induction.}
Delete a leaf $u$ from a tree $T$ and suppose the remaining tree has
been embedded. If $u$ is an out-leaf at its parent $v$, extension
requires an unused outneighbour of the image of $v$; for an in-leaf
it requires an unused inneighbour. The average-degree calculation
above finds a suitable centre somewhere in the tournament, but the
image of $v$ in an already fixed embedding need not be that vertex.
It may have used all neighbours of the required orientation.

For a minimal illustration of this logical failure, a transitive
three-vertex tournament has a sink with no outneighbour. An embedding
of a one-vertex parent at that sink cannot be extended to an out-edge,
although an out-edge certainly embeds elsewhere. Thus a failed greedy
extension is not a counterexample, and counting the globally unused
vertices alone does not repair it. A successful induction needs control
of the possible parent images, or a rerouting theorem preserving the
already embedded branches. No such uniform rerouting lemma is proved
in this attempt.

\paragraph{Verification record and exact remaining gap.}
The tournament enumeration is in \texttt{checks\_graphs.py} and
\texttt{results\_graphs.json} under
\path{_Local/Erdos_problems/UnsolvedMath/attempt_audit_2026_10_01/number_theory/}.
The sharp lower bound holds for every $n$, and the sharp upper bound
is proved here for the two star orientations and for $n\le3$, with
an even stronger result for consistently directed paths. Arbitrary
branching orientations and the full uniform finite range remain
uncovered by these deductions. The large-order theorem is reported
with its actual quantifier, not converted into an unperformed finite
verification. \textbf{Outcome: UNRESOLVED.}

\subsection{Sources}

{\small

\begin{itemize}

\item[\hypertarget{source:1424:1}{S1}] ULAM.AI, \emph{UnsolvedMath}, supplied \texttt{problems.json}, record 1424 (GRAPH-037), statement and background. The embedded literature-review date is 2026-08-17; that is the archive’s review date, not a fresh certification. Dataset landing page: \url{https://huggingface.co/datasets/ulamai/UnsolvedMath}.

\item[\hypertarget{source:1424:2}{S2}] Daniela Kuhn, Richard Mycroft, and Deryk Osthus, A proof of Sumner's universal tournament conjecture for large tournaments, Proceedings of the London Mathematical Society 102 (2011), 731-766. \url{https://arxiv.org/abs/1010.4430}. Bibliographic information imported from the supplied record.

\item[\hypertarget{source:1424:3}{S3}] Jiangdong Ai and Xiaopan Lian, An improved finite bound for oriented trees in tournaments, arXiv:2608.11667 (2026). \url{https://arxiv.org/abs/2608.11667}. Bibliographic information imported from the supplied record.


\item[E1] Daniela K\"uhn, Richard Mycroft, and Deryk
Osthus, \emph{A proof of Sumner's universal tournament conjecture for
large tournaments}, arXiv:1010.4430; Proc. London Math. Soc. 102 (2011),
731--766. \url{https://arxiv.org/abs/1010.4430}.
Primary abstract inspected 1 October 2026.
\item[E2] Jiangdong Ai and Xiaopan Lian, \emph{An improved finite
bound for oriented trees in tournaments}, arXiv:2608.11667v1.
\url{https://arxiv.org/abs/2608.11667}. Primary abstract inspected
1 October 2026; the uniform bound reported there is distinguished from
the conjectured sharp bound.

\end{itemize}

}

\label{end:1424}
\end{document}
